% !TEX root = ../Main.tex
\chapter{空间直角坐标系与空间向量}

\section{空间直角坐标系（右手系）}

在空间取三条两两垂直、公共原点 $O$ 的数轴 $x,y,z$，方向按\textbf{右手系}规定。空间中任一点都对应唯一一组坐标 $(x,y,z)$。

\begin{center}
\begin{tikzpicture}[>=Stealth, scale=1]
    \draw[->, thick] (0,0) -- (0,2.4) node[above] {$z$};
    \draw[->, thick] (0,0) -- (2.6,0.6) node[right] {$y$};
    \draw[->, thick] (0,0) -- (-1.3,-1.3) node[below left] {$x$};
    \node[below right] at (0,0) {$O$};
\end{tikzpicture}
\end{center}

\section{空间向量基本定理}

\state{空间向量基本定理}{
    $\vec e_1,\vec e_2,\vec e_3$ 不共面，则空间中的任意一个向量 $\vec a$，都可以表示为
    \[
    \vec a=x\vec e_1+y\vec e_2+z\vec e_3.
    \]
}

“不共面”是空间版本的关键条件——三条不共面的向量才撑得起整个空间；若三者共面，就只能张成一个平面。

\state{推论：四点共面}{
    设 $A,B,C,D$ 四点共面，在外面任取一点 $O$ 当参考点，则
    \[
    \vec{OA}=m\,\vec{OB}+n\,\vec{OC}+p\,\vec{OD},\qquad m+n+p=1.
    \]
}

\section{坐标、两点向量、单位向量}

特别地，当 $|\vec e_1|=|\vec e_2|=|\vec e_3|=1$ 且 $\vec e_1,\vec e_2,\vec e_3$ 两两互相垂直时，
\[
\vec a=x\vec e_1+y\vec e_2+z\vec e_3\ \Longrightarrow\ \vec a=(x,y,z).
\]
这时点乘、模长、夹角与平面情形完全同形，只是各多一项 $z$：
\[
\vec a\cdot\vec b=x_1x_2+y_1y_2+z_1z_2,\qquad |\vec a|=\sqrt{x^2+y^2+z^2},\qquad
\cos\langle\vec a,\vec b\rangle=\frac{\vec a\cdot\vec b}{|\vec a|\,|\vec b|}.
\]

\textbf{两点向量}：$A(x_1,y_1,z_1),\ B(x_2,y_2,z_2)$，则
\[
\vec{AB}=(x_2-x_1,\ y_2-y_1,\ z_2-z_1).
\]

\textbf{方向单位向量的求法}：把非零向量 $\vec a$ 除以自己的模长，
\[
\vec a_0=\pm\frac{\vec a}{|\vec a|}.
\]
它表示 $\vec a$ 的方向、模长为 $1$；正负两支对应两个相反方向。
